% AI Development Log -- CE-402 Group 8
% Compile with pdfLaTeX (e.g. Overleaf: upload this file, Compiler = pdfLaTeX).
\documentclass[11pt,a4paper]{article}

\usepackage[utf8]{inputenc}
\usepackage[T1]{fontenc}
\usepackage{lmodern}
\usepackage[a4paper,margin=20mm,bottom=22mm]{geometry}
\usepackage{amsmath,amssymb}
\usepackage{xcolor}
\usepackage{array,tabularx,booktabs,longtable}
\usepackage{enumitem}
\usepackage{fancyhdr}
\usepackage{titlesec}
\usepackage[most]{tcolorbox}
\usepackage[colorlinks=true,linkcolor=blue!60!black,urlcolor=blue!60!black]{hyperref}

\definecolor{navy}{HTML}{0F172A}
\definecolor{accent}{HTML}{1D4ED8}
\definecolor{softbg}{HTML}{F8FAFC}
\definecolor{rule}{HTML}{93C5FD}

\setlength{\parindent}{0pt}
\setlength{\parskip}{5pt}
\setlist{nosep,leftmargin=*}

\titleformat{\section}{\Large\bfseries\color{navy}}{}{0pt}{}[\vspace{-4pt}\color{black!25}\rule{\textwidth}{0.6pt}]
\titleformat{\subsection}{\large\bfseries\color{accent}}{}{0pt}{}

\pagestyle{fancy}
\fancyhf{}
\renewcommand{\headrulewidth}{0pt}
\fancyfoot[L]{\footnotesize\color{black!60}CE-402 Geotechnical Design Platform (Group 8) --- AI Development Log}
\fancyfoot[R]{\footnotesize\color{black!60}Page \thepage}

\newtcolorbox{prompt}{breakable,enhanced,colback=softbg,colframe=rule,boxrule=0pt,leftrule=2.5pt,
  arc=0pt,left=8pt,right=6pt,top=4pt,bottom=4pt,fontupper=\small}
\newcommand{\stamp}[1]{{\small\bfseries\color{black!70}\textcolor{accent}{\rule[-2pt]{2pt}{10pt}}\hspace{6pt}#1}\par}
\newcommand{\code}[1]{\texttt{\small #1}}

\begin{document}

{\small\bfseries\color{accent}CE-402 GEOTECHNICAL ENGINEERING --- GROUP 8}\par
{\LARGE\bfseries AI Development Log}\par
{\color{black!65}Geotechnical Design Platform (Modules P1--P5 \& P8) and GeoFound --- verbatim prompts, build notes and design reasoning}\par
\vspace{2pt}{\color{navy}\rule{\textwidth}{1.5pt}}

How AI tools were used to build the two versions of our platform. Prompts are reproduced verbatim (long prompts in
full, follow-ups as written). The geotechnical content --- methods, equations, parameters, design decisions and the
checks --- was specified and verified by the group; the AI was used to write code to that specification. Every
calculation error we found in AI output is recorded in the companion document \emph{Engineering Audit of AI Output}.

\begin{center}
\begin{tabularx}{\textwidth}{>{\raggedright}p{0.48\textwidth}lX}
\toprule
\textbf{Version} & \textbf{Tool} & \textbf{Where} \\
\midrule
v1 --- GeoTech Analytics (P1--P5 engines, vanilla JS) & Gemini & \code{platform-v1/} \\
v2 --- GeoFound (web app: UI/UX, integration, design workflow) & Claude Code (VS Code) & repository root \\
\bottomrule
\end{tabularx}
\end{center}

%==========================================================================
\section{Part A --- v1 platform (Gemini)}

\subsection{Prompt 1 --- Foundational architecture (\code{index.html}, \code{js/state.js})}
\stamp{20.11.2026 \textbar\ 23:24}
\begin{prompt}
You are a senior geotechnical software engineer. We are building a browser-based, zero-installation geotechnical
engineering web platform (CE-402, Group 8) using HTML5, Tailwind CSS (via CDN), Plotly.js (via CDN), and modular
Vanilla JavaScript for direct hosting on GitHub Pages with no build tools or backend servers.

In this first step, generate ONLY the foundational architecture across two files:

\textbf{index.html:}
\begin{itemize}
\item Include CDN links for Tailwind CSS and Plotly.js.
\item Build a clean, modern engineering dashboard layout (slate-navy header bar showing `CE-402 Geotechnical Design
  Platform \textbar\ Group 8', responsive container, and clean tab navigation buttons for Module P1: Soil Profile,
  Module P2: SPT \& Direct Shear, Module P3, Module P4, and Module P5).
\item Include empty container \code{<div>} panels for each tab (\code{\#tab-p1} through \code{\#tab-p5}) and a small
  tab-switching script so clicking any tab cleanly toggles the active view (with \code{\#tab-p1} active by default).
\item Load standard \code{<script>} tags at the bottom (do NOT use \code{type="module"} so the app works both by
  opening \code{index.html} locally and on GitHub Pages) in this order: \code{js/state.js},
  \code{js/moduleP1\_logic.js}, \code{js/moduleP1\_ui.js}.
\end{itemize}

\textbf{js/state.js:}
\begin{itemize}
\item Create a global reactive state store \code{window.GeotechState} holding: \code{gamma\_w = 9.81} (kN/m$^3$),
  \code{zw\_measured} (default 3.50 m), \code{zw\_design} (default 1.50 m), raw \code{layers} array (pre-populated
  with 3 realistic default layers: Sand SW/SP 0--3 m, Clay CL/CH 3--9 m, Gravel GW/GP 9--14 m using all standard SI
  geotechnical parameters), \code{processedLayers}, and \code{validationErrors}.
\item Include \code{subscribe(callback)} and \code{notify()} methods so any future module (P2--P5) automatically
  recalculates whenever the user edits the soil profile or groundwater depths.
\end{itemize}
Do not build the Module P1 calculations or UI yet---just deliver complete, clean code for \code{index.html} and
\code{js/state.js}.
\end{prompt}

\subsection{Prompt 2 --- Module P1 calculation and validation engine (\code{js/moduleP1\_logic.js})}
\stamp{21.11.2026 \textbar\ 00:47}
\begin{prompt}
Now build Step 2: create \code{js/moduleP1\_logic.js}, attached to \code{window.ModuleP1Logic} (classic script using
an IIFE with \code{"use strict"}, no ES6 import/export so it runs on \code{file://} and GitHub Pages).

This file must contain ONLY the pure geotechnical validation and calculation engine for Module P1, strictly enforcing
SI units (kN, kPa, m, mm, kN/m$^3$) with \code{gamma\_w = 9.81} kN/m$^3$ (defaulting to
\code{window.GeotechState.gamma\_w} if available, otherwise 9.81).

Match the exact input layer object property names already defined in \code{js/state.js}: \code{soilType},
\code{z\_top}, \code{z\_bot}, \code{gamma\_moist}, \code{gamma\_sat}, \code{su}, \code{c\_prime}, \code{phi\_prime},
\code{E\_prime}, \code{nu\_prime}, \code{e0}, \code{Cc}, \code{Cr}, \code{OCR}, \code{cv}.

Implement a function \code{validateAndProcessProfile(layers, gammaW)} that does NOT mutate the input array, sorts a
shallow copy of the layers by \code{z\_top} ascending, and returns \code{\{ processedLayers, errors \}}:

\textbf{Mandatory Validation Checks} (append clear, layer-indexed error messages to \code{errors} if violated):
\begin{itemize}
\item At least one soil layer must exist.
\item Every required numeric field must be a finite number (no NaN, null, undefined, or Infinity).
\item First layer must start at ground surface ($|z_\text{top} - 0.0| \le 10^{-3}$ m).
\item Positive thickness for every layer ($H_m = z_\text{bot} - z_\text{top} > 0$).
\item Stratigraphic continuity between adjacent layers ($|z_\text{top}[i] - z_\text{bot}[i-1]| \le 10^{-3}$ m, no
  gaps or overlaps).
\item Phase relationship validity (\code{gamma\_moist > 0} and \code{gamma\_sat >= gamma\_moist}).
\item Elastic stability for drained Poisson's ratio ($0.0 \le \nu' < 0.499$) to prevent division-by-zero singularity
  in $E_\text{oed}$.
\item Physical bounds: $s_u \ge 0$, $c' \ge 0$, $\phi' \ge 0$ (and $\phi' \le 60$), $E' > 0$, $e_0 > 0$,
  $C_c > 0$, $C_r > 0$, $\text{OCR} \ge 1.0$, $c_v > 0$.
\end{itemize}

\textbf{Derived Engineering Calculations} (compute for every layer when its inputs are numerically valid, and attach
these exact property names to each object in \code{processedLayers} alongside the original fields):
\begin{itemize}
\item \code{H\_m} $= z_\text{bot} - z_\text{top}$ (m); \quad \code{H\_mm} $= 1000\,H_m$ (mm)
\item \code{gamma\_sub} $= \gamma_\text{sat} - \gamma_w$ (kN/m$^3$)
\item \code{E\_oed} $= \dfrac{E'(1-\nu')}{(1+\nu')(1-2\nu')}$ (kPa); \quad \code{mv} $= 1/E_\text{oed}$ (1/kPa)
\end{itemize}

\textbf{Also expose on \code{window.ModuleP1Logic}:}
\begin{itemize}
\item \code{SOIL\_STYLES}: a dictionary mapping the 6 USCS categories (``Clay (CL/CH)'', ``Silt (ML/MH)'', ``Sand
  (SW/SP)'', ``Gravel (GW/GP)'', ``Peat (Pt)'', ``Rock / Bedrock'') to clean, muted geological fill colors, border
  colors, and short codes for Plotly rendering.
\item \code{SOIL\_TYPES}: an ordered array of those exact 6 USCS category strings for dropdown menus.
\end{itemize}
Include concise engineering comments tracing every equation and validation rule to standard references (Terzaghi,
Peck \& Mesri; Bowles; Das \& Sobhan; Eurocode 7 EN 1997-1; ASTM D2487). Deliver ONLY the complete
\code{js/moduleP1\_logic.js} file.
\end{prompt}

\subsection{Prompt 3 --- Module P1 interactive UI (\code{js/moduleP1\_ui.js})}
\stamp{21.11.2026 \textbar\ 01:01}
\begin{prompt}
Now build Step 3 (final step for Module P1): create \code{js/moduleP1\_ui.js}, attached to \code{window.ModuleP1UI}
(classic IIFE with \code{"use strict"}, no ES6 import/export so it runs cleanly on \code{file://} and GitHub Pages).
This file renders the complete interactive Module P1 interface inside \code{\#tab-p1} of \code{index.html} using
\code{window.GeotechState} and \code{window.ModuleP1Logic} (including \code{SOIL\_TYPES}, \code{SOIL\_STYLES}, and
\code{validateAndProcessProfile}).

\textbf{Required UI \& Functional Specifications inside \code{\#tab-p1}:}
\begin{enumerate}
\item \textbf{Top KPI Summary Bar} (4 clean Tailwind cards): Total Profile Depth (m), Number of Layers, Measured Water
  Table \code{zw\_measured} (m), and Design Water Table \code{zw\_design} (m).
\item \textbf{Groundwater \& Profile Controls Bar:} Editable numeric inputs for \code{zw\_measured} (m) and
  \code{zw\_design} (m) (min 0.0, step 0.1) bound to \code{window.GeotechState.setWaterTables()}; a clean badge
  displaying $\gamma_w = 9.81$ kN/m$^3$; and a ``Reset to Sample Profile'' button wired to
  \code{window.GeotechState.resetToDefaults()}.
\item \textbf{Editable Stratigraphic Input Table} (no hard-coded calculations; must accept any unseen soil profile):
  Responsive horizontally scrollable Tailwind table covering all 15 parameters matching \code{js/state.js}: Soil Type
  (dropdown populated from \code{window.ModuleP1Logic.SOIL\_TYPES}), \code{z\_top}, \code{z\_bot},
  \code{gamma\_moist}, \code{gamma\_sat}, \code{su}, \code{c\_prime}, \code{phi\_prime}, \code{E\_prime},
  \code{nu\_prime}, \code{e0}, \code{Cc}, \code{Cr}, \code{OCR}, \code{cv}. Include an ``Add Layer'' button
  (automatically setting the new layer's \code{z\_top} to the last layer's \code{z\_bot} and \code{z\_bot} to
  \code{z\_top + 3.0 m}, with valid default Sand parameters) and a ``Delete'' button on each row.

  \emph{IMPORTANT UX DETAIL:} Do NOT destroy and re-render the entire input table DOM on every keystroke while the user
  is typing in a numeric cell (which causes focus loss). Update only the KPI cards, validation box, derived table, and
  Plotly chart on input/change events, and re-render the input table only on initial load, Add Layer, Delete Layer,
  or Reset.
\item \textbf{Validation Alert Banner:} Run
  \code{window.ModuleP1Logic.validateAndProcessProfile(state.layers, state.gamma\_w)} on every update and push the
  results to \code{window.GeotechState.setProcessedProfile(processedLayers, errors, true)} (ensure no infinite loop if
  subscribed). If \code{errors.length > 0}, display a prominent red Tailwind alert box listing every validation error
  and hide the Derived Parameters Table and Plotly chart until fixed.
\item \textbf{Two-Column Results Workspace} (shown only when \code{errors.length === 0}):
  \begin{itemize}
  \item \emph{Left Column:} Clean Derived Parameters Table displaying Layer \#, Soil Type (with small color swatch
    from \code{SOIL\_STYLES}, falling back safely if unknown), Depth Interval (\code{z\_top--z\_bot} m), \code{H\_m}
    (m), \code{H\_mm} (mm), \code{gamma\_sub} (kN/m$^3$), \code{E\_oed} (kPa), and \code{mv} (1/kPa in scientific
    notation, e.g., \code{toExponential(3)}).
  \item \emph{Right Column:} Interactive Plotly.js 2D Stratigraphic Borehole Column (\code{\#p1-profile-plot}) with
    downward-positive depth axis (\code{autorange: "reversed"}) labeled ``Depth below Ground Surface, z (m)'', cleanly
    padded slightly below the maximum profile depth; scaled rectangular lithology blocks colored with
    \code{SOIL\_STYLES} \code{fillColor} and \code{borderColor}, with centered text annotations showing USCS code,
    soil type, and depth/thickness; invisible scatter/marker hover targets centered in each layer (or polygon hover
    traces) displaying unit weights, strength ($s_u$, $c'$, $\phi'$), stiffness ($E'$, $\nu'$, $E_\text{oed}$,
    $m_v$), and consolidation parameters ($e_0$, $C_c$, $C_r$, OCR, $c_v$); distinct horizontal lines and right-side
    annotations for Measured WT (dashed blue) and Design WT (solid red); using
    \code{Plotly.react(..., \{ responsive: true, displaylogo: false \})} for smooth updates.
  \end{itemize}
\end{enumerate}
Initialize automatically on \code{DOMContentLoaded}. Deliver ONLY the complete \code{js/moduleP1\_ui.js} file.
\end{prompt}

\subsection{Prompt to improve the solution to Prompt 3}
\stamp{21.11.2026 \textbar\ 01:21}
\begin{prompt}
In \code{js/moduleP1\_ui.js}, fix the following execution-order bug and apply one small Plotly UX improvement, then
output the complete updated \code{js/moduleP1\_ui.js} file:
\begin{enumerate}
\item \textbf{Bug Fix (Critical):} In both \code{handleAddLayer()} and \code{handleDeleteLayer()},
  \code{renderInputTable()} is currently called BEFORE \code{state.setLayers(nextLayers)}. Because
  \code{renderInputTable()} reads from \code{state.layers}, the editable table stays out of sync when adding or
  deleting rows. Update both functions so \code{state.setLayers(nextLayers)} is called BEFORE
  \code{renderInputTable()}.
\item \textbf{Minor Cleanup:} In the \code{btnReset} click handler, remove the redundant extra call to
  \code{refreshDerivedAndVisuals(state)} since \code{state.resetToDefaults()} already notifies subscribers (just keep
  \code{state.resetToDefaults()} followed by \code{renderInputTable()}).
\item \textbf{Plotly UX Improvement:} Instead of a single small invisible hover marker at the center of each layer,
  render an invisible filled polygon trace (\code{type: "scatter", mode: "lines", fill: "toself", hoveron: "fills",
  fillcolor: "rgba(0,0,0,0)", line: \{ width: 0 \}}) spanning each layer rectangle \code{[colX0, colX1, colX1, colX0,
  colX0]} by \code{[z\_top, z\_top, z\_bot, z\_bot, z\_top]} so hovering anywhere inside a soil stratum displays its
  full parameter tooltip.
\end{enumerate}
Keep all other logic, Tailwind classes, IDs, SI units, and classic IIFE structure unchanged. Deliver ONLY the complete
updated \code{js/moduleP1\_ui.js} file.
\end{prompt}

\subsection{Phase 2 --- Module P2: SPT corrections and direct-shear parameter fitting}
\stamp{21.11.2026 \textbar\ 22:01 (Step 2A) \quad 22.11.2026 \textbar\ 17:50--19:20 (Step 2C)}
\begin{prompt}
We are now moving to Phase 2 (Module P2: SPT Corrections \& Direct Shear Parameter Fitting). For Step 2A, update ONLY
\code{js/state.js} (classic IIFE with \code{"use strict"}, no ES6 modules, attached to \code{window.GeotechState}).

CRITICAL CONSTRAINT: Preserve 100\% of the existing Module P1 properties, default layers, and methods
(\code{gamma\_w}, \code{zw\_measured}, \code{zw\_design}, \code{layers}, \code{processedLayers},
\code{validationErrors}, \code{subscribe}, \code{notify}, \code{setWaterTables}, \code{setLayers},
\code{setProcessedProfile}, \code{resetToDefaults}) with exact backward compatibility so Module P1 continues working
unchanged.

Extend \code{window.GeotechState} with a Module P2 state slice using strict SI units (m, kPa, kN/m$^3$):
\begin{enumerate}
\item \code{sptRecords}: pre-populated via \code{createDefaultSptRecords()} with 5 realistic editable starter SPT
  observations spanning the default 0--14 m profile, each using these exact property names: \code{depth\_m} (number,
  e.g., 1.5, 3.0, 4.5, 7.5, 10.5), \code{N\_measured} (number, e.g., 12, 18, 10, 14, 28), \code{soilCategory}
  (string, one of: ``Fine Sand / Silt (Dilatancy Applicable)'', ``Coarse Sand / Gravel (No Dilatancy)'', ``Cohesive
  Soil (Clay)''), \code{sigma\_v\_override} (null by default, meaning auto-compute effective overburden pressure from
  Module P1 stratigraphy \& \code{zw\_measured}; or a positive number in kPa if manually overridden).
\item \code{directShearPoints}: pre-populated via \code{createDefaultDirectShearPoints()} with 4 realistic editable
  starter Direct Shear test observations, each using these exact property names: \code{sigma\_n} (number, normal
  stress in kPa, e.g., 50, 100, 150, 200), \code{tau\_f} (number, peak shear stress in kPa, e.g., 38.5, 69.0, 96.5,
  128.0).
\item \code{p2Results}: \code{\{ sptProcessed: [], sptErrors: [], directShearFit: null, directShearErrors: [] \}}
\item Add clean mutators on \code{window.GeotechState}: \code{setSptRecords(records)} (updates \code{sptRecords} and
  calls \code{notify()}), \code{setDirectShearPoints(points)} (updates \code{directShearPoints} and calls
  \code{notify()}), \code{setP2Results(resultsObj, triggerNotify)} (stores \code{p2Results} without auto-looping
  unless \code{triggerNotify} is true), and \code{resetP2Defaults()} (resets only \code{sptRecords},
  \code{directShearPoints}, and \code{p2Results} back to defaults and calls \code{notify()}).
\end{enumerate}
Keep \code{resetToDefaults()} resetting P1 as before (and optionally leave P2 intact so resetting the soil profile does
not erase lab test inputs). Deliver ONLY the complete updated \code{js/state.js} file.
\end{prompt}

\textbf{What was built.} Step 2A extended the state store for P2 with full backward compatibility. Step 2B
(\code{js/moduleP2\_logic.js}) implemented piecewise effective overburden split at the measured water table, the
IS 2131:1981 overburden correction ($\sigma'_v \ge 25$ kPa, $C_N \le 2.0$) and dilatancy correction ($N' > 15$ below
the water table in fine sands/silts) with step-by-step traces, and ordinary least-squares Mohr--Coulomb fitting ($c'$,
$\phi'$, $R^2$, RMSE, residuals) with a non-negative cohesion check. Step 2C (\code{js/moduleP2\_ui.js}) added the
editable SPT and direct-shear tables, calculation cards and Plotly charts. Error found in review: audit item A1.

\subsection{Phase 3 --- Module P3: IS 6403:1981 bearing capacity}
\stamp{22.11.2026 \textbar\ 21:11--23:00 \quad audit 26.11.2026 \textbar\ 01:01--02:11}
Step 3A extended the state store (17 strict-SI geometry, loading, shear-mode and soil-source fields). Step 3B
(\code{js/moduleP3\_logic.js}) implemented input and eccentricity validation, Meyerhof (1953) effective dimensions for
strip, square, rectangular and circular footings, thickness-weighted failure-zone strength from the P1 profile,
IS 6403 shear modes and interpolation, shape, depth, inclination and $W'$ factors, net and gross ultimate and safe
capacities, audit trails and a footing-width sensitivity curve. Step 3C (\code{js/moduleP3\_ui.js}) added the KPI
cards, factor table, calculation log and charts. Error found in review: audit item A2.

\subsection{Phase 4 --- Modules P4/P5: settlement (IS 8009 Part 1:1976, IS 1904:1986)}
\stamp{26.11.2026 \textbar\ Step 4C 05:52--07:30}
Step 4A extended the state store (\code{p4Input}, \code{p4Results}, \code{deepMerge}/\code{cloneValue},
\code{syncP4FromP3}). Step 4B (\code{js/moduleP4\_logic.js}) was directed to implement 2:1, Newmark/Boussinesq,
circular, strip and Westergaard stress distribution; IS 8009/Janbu elastic settlement with the Fox depth correction
and Schmertmann (1978) strain influence; 1-D consolidation with the three NC/OC cases, a 25-term Fourier $U(T_v)$,
drainage path, $t_{50}/t_{90}$ and time-rate curves. Step 4C (\code{js/moduleP4\_ui.js}) added the settlement UI with
IS 1904 presets and three charts. Errors found in review: audit items A3--A7.

\subsection{Step 5B and P8 --- stress-increase module and deployment}
Exact Newmark (1935)/Boussinesq vertical stress below rectangular corners and centres by four-quadrant superposition,
2:1 dispersion, depth-profile chart and audit trail (\code{js/moduleP4\_stress\_logic.js},
\code{js/moduleP4\_stress\_ui.js}); README added (P8). Error found in review: audit item A8.

%==========================================================================
\section{Part B --- v2 GeoFound: UI/UX (Claude Code)}
\stamp{26.11.2026}
GeoFound is the web app that runs the design end-to-end. This part logs the UI/UX work only; the engine, the
foundation design and every calculation decision are covered in the engineering audit and in
\code{tests/hand\_calcs.md}.

\subsection{Prompts (verbatim)}
\begin{enumerate}[label=\arabic*.]
\item ``hey i want really good design that's why i chose figma and i want it dynamic''
\item ``not that good design tell me any connector i should install for the reference of the ui ux''
\item ``i am giving the site links for the reference design the site whose ui should be used as a reference
  \url{https://busybox.io/} and \url{https://www.neuroformtech.com/} i have provided two references either you can use
  any one of these or make a new one using this reference''
\item ``ui looks good but any different colour can you check for another colour as well of the site i mean it's kinda
  dark brown''
\item ``go with the dark brown colour it looks good''
\item ``one more thing pick some components for boxes and fonts cause it looks like ai generated and the font and box
  is not looking good cause it's edges are sharp and could be better if it hover when we place cursor over it you know
  and font is not that good i mean if you can find the better font style then use that one the box and all cause it
  still looks like ai''
\item ``can you take reference from \url{https://www.neuroformtech.com/} for dynamic site like it's good and continue
  further now''
\item ``isn't while moving down texts are overlapping and i want some more dynamic component like adding video in the
  background which will be moving continuously and all''
\item ``can you add that moving part which you added in front one for all or something dynamic in the background which
  moves continuously to make it dynamic''
\item ``\ldots\ when i am clicking it is not going to home page where geofound is writtern and soil profile is shown it
  directly take us to design page but i want it to go to home page and then person can go anywhere also it get bit
  unresponsive sometimes and i have to reload it \ldots''
\item ``\ldots\ i am thinking that for the new site data whether a user can enter it or upload the csv file for that''
\item ``\ldots\ when i click on this soil profile then it detaches and show it separately like layerwise are you
  getting na''
\item ``use ony free api key no claude one cause it's paid''
\item ``can a person gives the input about seismic zone ??''
\end{enumerate}

\subsection{AI reasoning (summarised from the session)}
\begin{itemize}[leftmargin=12pt]
\item \textbf{Design direction.} Four directions were tried. The dark-brown ``Core Sample'' direction reads as a soil
  core rather than a generic dashboard: ground \#140806, warm ink \#F4E8DF, and survey-marking orange \#FF4A1C kept for
  actions and pass/fail only, so colour still carries meaning on result screens.
\item \textbf{Why it looked AI-made.} Inter / IBM Plex and sharp 1 px boxes are the default look of generated
  dashboards. They were replaced with a display face for headlines (Tanker), a grotesk for text (Cabinet Grotesk,
  Satoshi) and a monospace with tabular figures for numbers (JetBrains Mono), so columns of results line up. Cards got
  rounded corners, a cursor-following spotlight and a small lift on hover.
\item \textbf{Background video.} A 5--15 MB video loop would cost load time and battery on every page. A WebGL shader
  that draws slowly moving strata gives continuous motion with zero download; it is capped at 30 fps behind content
  pages and switched off under \code{prefers-reduced-motion}.
\item \textbf{Text overlapping on scroll.} The landing page's giant word used scroll parallax that moved it up, over
  the buttons. It now sinks and fades behind the ticker, and the navigation bar has a solid background.
\item \textbf{Navigation ``jumping to Design'' and freezing.} The URL changed but the old page stayed on screen:
  \code{AnimatePresence mode="wait"} kept the previous route mounted until its exit animation finished, and rapid
  clicks during a heavy exit left it stuck. The wait mode was removed and the fix was checked with rapid clicks, the
  logo link and Back.
\item \textbf{Exploding soil core.} Framer Motion's \code{y} on an SVG \code{<g>} does not move the group, so the layers
  did not separate. A CSS \code{transform: translateY} with a transition is used, and each layer card is placed below
  the previous one so labels never overlap.
\item \textbf{CSV upload.} One CSV with three record types --- borehole, layer, spt --- so a site can be prepared in a
  spreadsheet; a downloadable template and a report of what was imported.
\item \textbf{Seismic zone input.} A zone dropdown (II--V, IS 1893 Table 3) is what a user thinks in; it fills $Z$ and
  $a_\text{max} = Z/2$ and still allows a custom value. Every hard-coded ``zone V'' label was removed.
\end{itemize}

\subsection{Accepted}
\begin{itemize}[leftmargin=12pt]
\item Visual direction ``Core Sample'' (dark brown + survey orange), chosen over three other directions and three
  colour variants; fonts Tanker, Cabinet Grotesk, Satoshi, JetBrains Mono; rounded hover cards.
\item Motion: loader, liquid headline, scramble hover, page transitions, live WebGL strata on every page, tickers,
  magnetic buttons, animated charts --- all off under \code{prefers-reduced-motion}.
\item Exploding soil core on the landing page (click or Enter), built from the example borehole data.
\item Data entry: step-by-step wizard, CSV import with template, project save/open as JSON, seismic zone dropdown.
\item AI summary on the Summary page with the Google Gemini free tier, a deterministic template fallback, a number
  check that flags any value not in the results, and an ``Engineer checked this'' gate before AI text goes into the PDF.
\end{itemize}

\subsection{Changed}
\begin{itemize}[leftmargin=12pt]
\item Fonts and boxes (felt AI-generated) $\rightarrow$ new typefaces and rounded hover cards.
\item Landing parallax (text over buttons) $\rightarrow$ sinks and fades; solid navigation bar.
\item Route transitions (freeze) $\rightarrow$ wait mode removed.
\end{itemize}

\subsection{Rejected}
\begin{itemize}[leftmargin=12pt]
\item Paid AI API for the summary --- free tier only.
\item A real background video --- replaced by the GPU-drawn backdrop.
\item Completing every screen in Figma --- the Figma plan's tool-call limit was reached; the remaining screens were
  designed in code in the approved style.
\end{itemize}

\end{document}
